\subsection{Flat modules}

PROPOSITION 1. Let E be a right A-module. The following three properties are equivalent:

\begin{enumerate}
\def\labelenumi{(\alph{enumi})}
\item
  E is flat for \(A_{s}\) (in other words, for every left ideal a of A, the canonical homomorphism \(E \otimes_{A} a \to E \otimes_{A} A_{s} = E\) is injective).
\item
  E is M-flat for every left A-module M.
\item
  For every exact sequence of left A-modules and homomorphisms
\end{enumerate}

\[
\mathrm{M} ^ {\prime} \xrightarrow {v} \mathrm{M} \xrightarrow {w} \mathrm{M} ^ {\prime \prime}
\]

the sequence

\[
\mathrm{E} \otimes_ {\mathrm{A}} \mathrm{M} ^ {\prime} \xrightarrow {1 \otimes v} \mathrm{E} \otimes_ {\mathrm{A}} \mathrm{M} \xrightarrow {1 \otimes w} \mathrm{E} \otimes_ {\mathrm{A}} \mathrm{M} ^ {\prime \prime}
\]

is exact.

It is immediate that (b) implies (a). Conversely, suppose that (a) holds; by Lemma 5 of no. 2, E is flat for every free left A-module; as every left A-module is isomorphic to a quotient of a free module (Algebra, Chapter II, §1, no. 11, Proposition 20), it follows from Lemma 4 of no. 2 that E is flat for M.

We show that (c) implies (b). If \(v \colon \mathbf{M}' \to \mathbf{M}\) is an injective homomorphism, the sequence \(0 \to \mathbf{M}' \xrightarrow{v} \mathbf{M}\) is exact; by (c) the sequence

\[
0 \rightarrow \mathrm{E} \otimes_ {\mathrm{A}} \mathrm{M} ^ {\prime} \xrightarrow {1 \otimes v} \mathrm{E} \otimes_ {\mathrm{A}} \mathrm{M}
\]

is exact; this means that \(1 \otimes v\) is injective, in other words, E is M-flat.

Finally, the implication \((\mathbf{b}) \Rightarrow (\mathbf{c})\) is a consequence of the following more precise lemma:

LEMMA 6. If \(\mathbf{M}' \xrightarrow{v} \mathbf{M} \xrightarrow{w} \mathbf{M}''\) is an exact sequence of left A-modules and if E is an \(\mathbf{M}''\) -flat right A-module, the sequence

\[
\mathrm{E} \otimes_ {\mathrm{A}} \mathrm{M} ^ {\prime} \xrightarrow {1 \otimes v} \mathrm{E} \otimes_ {\mathrm{A}} \mathrm{M} \xrightarrow {1 \otimes w} \mathrm{E} \otimes_ {\mathrm{A}} \mathrm{M} ^ {\prime \prime}
\]

is exact.

We use the notation \(\mathrm{T}(\mathbf{Q})\) and \(\mathrm{T}(v)\) in the same sense as in the proof of Lemma 4 of no. 2. We write \(\mathbf{M}_1'' = w(\mathbf{M})\) and let \(i: \mathbf{M}_1'' \to \mathbf{M}''\) be the canonical injection and \(p\) the mapping of \(\mathbf{M}\) to \(\mathbf{M}_1''\) with the same graph as \(w\) . The sequence \(\mathbf{M}' \xrightarrow{v} \mathbf{M} \xrightarrow{p} \mathbf{M}_1'' \to 0\) is exact and it follows from Lemma 1 of no. 1 that the sequence

\[
\mathrm{T} (\mathbf {M} ^ {\prime}) \xrightarrow {\mathrm{T} (v)} \mathrm{T} (\mathbf {M}) \xrightarrow {\mathrm{T} (p)} \mathrm{T} (\mathbf {M} _ {1} ^ {\prime \prime}) \longrightarrow 0
\]

is exact. Moreover, as E is \(\mathbf{M}''\) -flat, the mapping \(\mathrm{T}(i): \mathrm{T}(\mathbf{M}_1'') \to \mathrm{T}(\mathbf{M}'')\) is injective, and as \(\mathrm{T}(i) \circ \mathrm{T}(p) = \mathrm{T}(w)\) , the sequence

\[
\mathrm{T} (\mathbf {M} ^ {\prime}) \xrightarrow {\mathrm{T} (v)} \mathrm{T} (\mathbf {M}) \xrightarrow {\mathrm{T} (w)} \mathrm{T} (\mathbf {M} ^ {\prime \prime})
\]

is exact. (§ 1, no. 3.)

DEFINITION 2. A right A-module E is called flat if it has the equivalent properties of Proposition 1.

Flat left A-modules are defined similarly. To say that a right A-module E is flat equivalent to saying that E, considered as a left A \(^{0}\) -module, is flat.

Remarks

\begin{enumerate}
\def\labelenumi{(\arabic{enumi})}
\item
  By Lemma 3 of no. 2, for a right A-module E to be flat, it is necessary and sufficient that, for every finitely generated left ideal a of A, the canonical mapping \(E \otimes_{A} a \to E\) (Proposition 1) with image Ea be injective.
\item
  Let E be a flat right A-module. If M' is a submodule of a left A-module M, the canonical injection \(E \otimes_{A} M' \to E \otimes_{A} M\) allows us to identify \(E \otimes_{A} M'\) with a subgroup of \(E \otimes_{A} M\) . This being so, let N be a left A-module, \(u: M \to N\) a homomorphism, K = Ker u, and I = Im u. By considering the exact sequence
\end{enumerate}

\[
0 \longrightarrow \mathrm{K} \longrightarrow \mathrm{M} \stackrel {u} {\longrightarrow} \mathrm{N}
\]

it is easily seen (Proposition 1) that \(\mathbf{E} \otimes_{\mathbf{A}} (\operatorname{Ker} u)\) is identified with \(\operatorname{Ker}(1_{\mathbf{E}} \otimes u)\) . On the other hand, writing \(u'\) for the surjective homomorphism \(\mathbf{M} \to \mathbf{I}\) with the same graph as \(u\) , and \(i\) for the canonical injection \(\mathbf{I} \to \mathbf{N}\) , \(1_{\mathbf{E}} \otimes u'\) is surjective (no. 1, Lemma 1) and \(1_{\mathbf{E}} \otimes i\) is injective since \(\mathbf{E}\) is flat. As

\[
\mathbf {1} _ {\mathbf {E}} \otimes \boldsymbol {u} = (\mathbf {1} _ {\mathbf {E}} \otimes i) \circ (\mathbf {1} _ {\mathbf {E}} \otimes u ^ {\prime}),
\]

\(E \otimes_{A} (Im u) \text{ is identified with } Im(1_{E} \otimes u).\)

PROPOSITION 2. (i) Let \((\mathbf{E}_t)_{t \in I}\) be a family of right A-modules. For \(\mathbf{E} = \bigoplus_{t \in I} \mathbf{E}_t\) to be flat, it is necessary and sufficient that each of the \(\mathbf{E}_t\) be flat.

\begin{enumerate}
\def\labelenumi{(\roman{enumi})}
\setcounter{enumi}{1}
\tightlist
\item
  Let \(\mathbf{I}\) be an ordered set and \((\mathrm{E}_{\alpha}, f_{\beta \alpha})\) a direct system of right A-modules (Algebra, Chapter II, § 6, no. 2). If each of the \(\mathrm{E}_{\alpha}\) is flat, then \(\mathbf{E} = \varinjlim \mathrm{E}_{\alpha}\) is flat.
\end{enumerate}

Let \(M' \to M\) be an injective left A-module homomorphism.

\begin{enumerate}
\def\labelenumi{(\roman{enumi})}
\tightlist
\item
  For the direct sum homomorphism
\end{enumerate}

\[
\bigoplus_ {i \in I} \left(E _ {i} \otimes_ {A} M ^ {\prime}\right)\rightarrow \bigoplus_ {i \in I} \left(E _ {i} \otimes_ {A} M\right)
\]

to be exact, it is necessary and sufficient that each of the homomorphisms \(\mathbf{E}_{\iota} \otimes_{\mathbf{A}} \mathbf{M}' \to \mathbf{E}_{\iota} \otimes_{\mathbf{A}} \mathbf{M}\) be so (Algebra, Chapter II, § 1, no. 6, Corollary 1 of Proposition 7), which proves (i) since \(\bigoplus_{\iota \in \mathbf{I}} (\mathbf{E}_{\iota} \otimes_{\mathbf{A}} \mathbf{M})\) is canonically identified with \(\mathbf{E} \otimes_{\mathbf{A}} \mathbf{M}\) (Algebra, Chapter II, § 3, no. 7, Proposition 7).

\begin{enumerate}
\def\labelenumi{(\roman{enumi})}
\setcounter{enumi}{1}
\tightlist
\item
  By hypothesis each of the sequences
\end{enumerate}

\[
0 \rightarrow \mathrm{E} _ {\alpha} \otimes_ {\mathrm{A}} \mathrm{M} ^ {\prime} \rightarrow \mathrm{E} _ {\alpha} \otimes_ {\mathrm{A}} \mathrm{M}
\]

is exact; so then is the sequence

\[
0 \rightarrow \mathrm{E} \otimes_ {\mathrm{A}} \mathrm{M} ^ {\prime} \rightarrow \mathrm{E} \otimes_ {\mathrm{A}} \mathrm{M}
\]

since taking the direct limit commutes with tensor products (Algebra, Chapter II, § 6, no. 3, Proposition 7) and preserves exactness (ibid., § 6, no. 2, Proposition 3).

